% Generated draft: requires mathematical and source review.
\documentclass[11pt]{article}
\usepackage[margin=25mm]{geometry}
\usepackage{fontspec}
\setmainfont{Latin Modern Roman}
\usepackage{amsmath,amssymb,mathtools,unicode-math}
\setmathfont{Latin Modern Math}
\usepackage{enumitem,needspace,xurl,hyperref}
\setlength{\parindent}{0pt}
\setlength{\parskip}{5pt}
\setlength{\emergencystretch}{3em}
\newcommand{\N}{\mathbb{N}}
\newcommand{\Z}{\mathbb{Z}}
\newcommand{\Q}{\mathbb{Q}}
\newcommand{\R}{\mathbb{R}}
\newcommand{\C}{\mathbb{C}}
\newcommand{\F}{\mathbb{F}}
\newcommand{\A}{\mathbb{A}}
\newcommand{\PP}{\mathbb{P}}
\newcommand{\E}{\mathbb{E}}
\newcommand{\eps}{\varepsilon}
\newcommand{\dd}{\,\mathrm{d}}
\begin{document}
\section*{UnsolvedMath ID: 20000930}
\textbf{Source problem code:} catalogue entry
\subsection{Definitions and mathematical statement}
A tournament is a finite directed graph with exactly one of $u\to v$ or $v\to u$ for each pair of distinct vertices. A subtournament is the induced tournament on a vertex subset; it is transitive if its vertices admit a linear ordering with every edge directed forward. Write $\operatorname{tr}(T)$ for the largest number of vertices in a transitive subtournament.

The original extract does not describe $T^*$ by its edges. The attached aid identifies it with the exceptional six-vertex tournament called $K_6$ by Berger--Choromanski--Chudnovsky, namely the Paley tournament on $\Z/7\Z$ with vertex $0$ removed. Under that supplied identification, use vertex set $\{1,2,3,4,5,6\}$ and orient $i\to j$ when $j-i\in\{1,2,4\}\pmod7$. The symbol $K_6$ here names this tournament, not an undirected complete-graph problem.

The workshop question is whether
\[
 \exists\epsilon>0\ \forall n\geq1\ \forall\text{ $n$-vertex tournaments }T:
 \quad T^*\not\subseteq T\ \Longrightarrow\ \operatorname{tr}(T)\geq n^\epsilon.
\]
Thus $\epsilon$ depends only on the fixed forbidden tournament, not on $n$ or $T$. The source's historical remark calls this the smallest then-unresolved tournament case, attributed to Berger--Choromanski--Chudnovsky. Its universal version quantifies separately for each forbidden tournament $H$: $\forall H\,\exists\epsilon_H>0$ with the corresponding bound for all $H$-free tournaments. The source states that this universal property is equivalent to the Erd\H{o}s--Hajnal conjecture; it does not ask for one exponent uniform over all $H$.

A tournament is two-colourable if its vertices partition into two transitive sets. For this $T^*$ the sets $\{1,2,3\}$ and $\{4,5,6\}$ are transitive. The supplied aid attributes to Nguyen, Scott, and Seymour (2025) the theorem that every two-colourable tournament has the stated Erd\H{o}s--Hajnal property, and reports that it answers this workshop question affirmatively. This is a supplied literature conclusion requiring independent verification, not a proof or source check performed here.

The aid also records a quantitative restricted class. Start with the one-vertex tournament and close under substitution into tournament templates with at most five vertices: replace each template vertex by a nonempty previously constructed block and orient every interblock edge according to the corresponding template edge. For every tournament $G$ in this recursive class, it claims
\[
 T^*\not\subseteq G,\qquad
 \operatorname{tr}(G)\geq |V(G)|^{\log_5 2}.
\]
No balance assumption on block sizes is made. This explicit exponent applies to the substitution class, not to all $T^*$-free tournaments as a consequence of the supplied summary.
\subsection{Short English statement}
Ask whether excluding the six-vertex Paley subtournament forces a transitive subtournament of polynomial size. The supplied aid attributes an affirmative answer to a 2025 theorem for two-colourable tournaments and gives an explicit exponent only for a recursively defined substitution class.
\subsection{Sources}
\begin{itemize}
\item AIM workshop \emph{Graph Ramsey theory}, section labeled ``Online Ramsey numbers,'' question 8, with historical remarks attributed to Berger--Choromanski--Chudnovsky: \url{https://www.overleaf.com/read/mnvcscjjysvg}.
\item The supplied aid identifies $T^*$ with their exceptional six-vertex tournament and attributes the two-colourable-tournament theorem to Nguyen, Scott, and Seymour (2025). No complete paper title or publication details are supplied.
\item Local registry \texttt{data/unsolvedmath/problems.json}, permanent ID 20000930, exact source code \texttt{catalogue entry}; canonical location \texttt{aim-combinatorics-notes.json notes[54]}. Attached attempt 1 is labeled full solution, problem status at run solved, rigor proved, and low novelty confidence. It is machine-generated, not peer reviewed, and requires independent verification before citation.
\end{itemize}
\end{document}
